{\parskip=0pt\centering\large\textbf{ABSTRACT}\par}

\addcontentsline{toc}{chapter}{ABSTRACT}
\vspace{3ex}

{\parskip=0pt\centering\textbf{\engtatitle{}}\par}
\vspace{\baselineskip}

\begingroup
\setlength{\tabcolsep}{0pt}

\bfseries
\begin{tabularx}{\textwidth}{l >{\centering}m{2em} X}
  Student Name / NRP & : & \name{} / \nrp{} \\[1ex]
  Department         & : & \engdepartment{} \engfacultyshort{} - ITS \\[1ex]
  Advisor            & : & Advisor Name and Title \\
\end{tabularx}
\endgroup
\vspace{\baselineskip}

\paragraphtextstyle

\noindent\textbf{Abstract}

% Ubah paragraf berikut dengan abstrak dari proposal tugas akhir dalam Bahasa Inggris
\lipsum[1]

\vspace{\baselineskip}

\noindent\textbf{Keywords:} \textbf{\textit{keyword 1, keyword 2, keyword 3, keyword 4, keyword 5}}.

